\section*{अध्याय \ref{ch:g1:five-senses} --- संसार को देखना: पाँच इंद्रियाँ}

\begin{solution}{exo:g1:five-senses:1}
दृष्टि --- आँखें; \omterm{def:g1:five-senses:senses}{श्रवण} --- कान; स्पर्श --- त्वचा; \omterm{def:g1:five-senses:senses}{घ्राण} --- नाक;
स्वाद --- जीभ।
\end{solution}

\begin{solution}{exo:g1:five-senses:2}
(क) \omterm{def:g1:five-senses:senses}{श्रवण}। (ख) स्वाद। (ग) स्पर्श। (घ) दृष्टि।
\end{solution}

\begin{solution}{exo:g1:five-senses:3}
केक को सबसे पहले \omterm{def:g1:five-senses:senses}{घ्राण} पकड़ता है --- महक दरवाज़े के नीचे से निकल आती
है। टुकड़ा खाने में दृष्टि, स्पर्श, \omterm{def:g1:five-senses:senses}{घ्राण} और स्वाद काम आते हैं, और
अगर ऊपरी परत कुरकुराती है तो \omterm{def:g1:five-senses:senses}{श्रवण} भी।
\end{solution}

\begin{solution}{exo:g1:five-senses:4}
उत्तर अलग-अलग होंगे; ज़्यादातर बच्चे $4$ या $5$ ध्वनियाँ गिनते हैं,
और चौंकाने वाली ध्वनि अकसर कोई धीमी-सी होती है --- फ़्रिज की भनभनाहट,
या अपनी ही साँस।
\end{solution}

\begin{solution}{exo:g1:five-senses:5}
जैसे: ``मेरा जूता नीला है, रबड़ का है, और मेरे हाथ से लंबा है।'' कोई
भी तीन सटीक शब्द सही हैं, बशर्ते कोई साथी उन्हें जाँच सके।
\end{solution}

\begin{solution}{exo:g1:five-senses:6}
अना का वाक्य \omterm{def:g1:five-senses:observation}{प्रेक्षण} है, जो स्पर्श से बना। बॉब का वाक्य राय है:
``सबसे अच्छी'' को देखकर या छूकर कोई जाँच नहीं सकता।
\end{solution}

\begin{solution}{exo:g1:five-senses:7}
(क) \omterm{def:g1:five-senses:senses}{घ्राण}। (ख) स्पर्श --- त्वचा थोड़ी दूर से भी गरमी महसूस कर लेती
है। (ग) \omterm{def:g1:five-senses:senses}{श्रवण}।
\end{solution}

\begin{solution}{exo:g1:five-senses:8}
अकेला स्पर्श तीनों के नाम आसानी से बता देता है: धातु का चम्मच (कठोर,
ठंडा, चिकना), स्पंज (मुलायम, हल्का), कंकड़ (कठोर, भारी, गोल)। चम्मच
आम तौर पर सबसे आसान रहता है --- ठंडी धातु को पहचानने में भूल कठिन है।
\end{solution}

\begin{solution}{exo:g1:five-senses:9}
पहली रेखा पर एक डोरी रखकर उसे उतनी ही लंबाई में काट लो या उँगली से
पकड़ लो, फिर उसे दूसरी रेखा पर रखो: डोरी दोनों की तुलना कर देती है,
और फ़ैसला आँखों को नहीं करना पड़ता। वही डोरी किसी टेढ़ी-मेढ़ी रेखा पर
भी चल सकती है और फिर सीधी खींची जा सकती है --- तो यह तरीक़ा टेढ़ी रेखा
की तुलना सीधी रेखा से भी कर देता है।
\end{solution}
